\section{Introduction}
\label{sec:intro}

Neural network weights are almost always written by an optimizer. A programmer chooses an architecture, a dataset and a loss, and gradient descent produces the numbers. On November 2, 2026, a blog post from MIT CSAIL reported that a frontier language model could do something different: read a network's weights and write or edit them directly, working from a natural-language specification without gradients \citep{takehi2026sparks}. The post called this \emph{weight-space programming} (WSP). In its experiments, the editor model named the game a reinforcement-learning policy was trained on from its weights alone. It repaired a car-control policy that never succeeded so that it succeeded in all 1{,}024 test episodes, by changing 3.5\% of the weights in one reply; a later audit found that the successful edits also work on a network with all weights set to zero, so they write a controller rather than repair the learned one \citep{zen24373915}. It also raised a 0.8B-parameter Qwen model's MATH-500 accuracy from 7.9\% to 48.1\% by editing 131{,}072 of its weights (0.016\%) \citep{takehi2026sparks}.

Within three weeks, dozens of follow-up studies appeared, on arXiv and in Zenodo preprints. They applied WSP to robot policies, clinical decision networks, quantum circuits, EEG and ECG models, materials science, tabular foundation models, vision-language models and diffusion models. They also replicated the original edits with other editor models. They found that a no-weight-change baseline (switching off the model's thinking mode) recovers most of the headline math gain \citep{turkoglu2026survey}, and that zero-shot weight reading is often local arithmetic plus priors rather than tracing computation through the network \citep{ax2611_08070}. Other studies showed that the same capability implants backdoors \citep{ax2611_17395,ax2611_27425,ax2611_27920}. The first studies of models editing their own weights appeared in the same window \citep{ax2611_21885,ax2611_28833,ax2611_29513}.

This survey organizes that literature. We cover \nIncluded\ studies (\nArxiv\ on arXiv and \nZenodo\ Zenodo preprints) posted up to \cutoffdate, each read in full or flagged where we could not. We code every study on four axes and place WSP inside the older work it draws on.

\paragraph{Why another survey now.} Two WSP surveys appeared during this period: a broad Chinese-language survey \citep{emecii2026survey} and an evidence audit with a reporting checklist \citep{turkoglu2026survey}. We build on both and do not repeat their contributions. In particular, we do not re-audit study-level evidence quality, and we refer readers to the WSP-12 checklist \citep{turkoglu2026survey} for reporting standards. \Cref{tab:surveys} compares the scope of all three with surveys of weight space learning, model merging and recursive self-improvement. None of the latter mentions an LLM acting as a weight editor. Both already give a definition, a taxonomy and open problems. Our survey differs from them in five specific ways.

\begin{enumerate}[leftmargin=*,itemsep=3pt]
\item \textbf{A definition and a taxonomy built for reuse} (\cref{sec:defs}). We give a one-sentence formal definition of WSP that separates it from hypernetworks, model editing and steering. We then classify every study on four axes: what the editor does (six branches, reading to self-editing), who authors the numbers (direct values to no LLM), the loop (one-shot, closed-loop or search), and the reference (another model or itself). The last three axes are aligned with \citet{turkoglu2026survey} so that the two surveys can be cross-referenced. Compared with the four axes of \citet{turkoglu2026survey}, our first axis adds feedback loops, efficient solvers and self-editing as branches, and our authorship axis adds a no-LLM level.
\item \textbf{Lineage} (\cref{sec:background}). We connect WSP to five older literatures, including fast weight programmers, which neither WSP survey covers: weights as programs (fast weight programmers, compiled transformers), weight space learning and its symmetries, hypernetworks, model editing, and self-improvement. Where those literatures have documented failure modes, such as symmetry sensitivity, ripple effects, forgetting and fragility under quantization, we say which ones WSP inherits.
\item \textbf{A reliability agenda} (\cref{sec:eval,sec:cost}). We go beyond auditing what papers report and ask what would have to be true for a WSP edit to be trusted. This covers reasoning versus pattern-matching, invariance under weight symmetries, formal and statistical verification, run-to-run variance, and whether edits survive fine-tuning and quantization. We find that persistence under fine-tuning and quantization is almost entirely untested. We also find that one reported cost figure for the original edit is not supported by the released logs.
\item \textbf{Conditions for recursive self-improvement (RSI)} (\cref{sec:rsi}). Current WSP is overwhelmingly one model editing another. We separate the early self-editing studies from claims of recursive self-improvement, and we state the conditions under which WSP could become a self-improvement loop.
\item \textbf{A numbered set of open problems} (\cref{sec:open}). Each problem comes with a concrete target that a skeptic can check, so future work can state which one it addresses.
\end{enumerate}

\begin{table}[t]
\centering
\footnotesize
\setlength{\tabcolsep}{3.5pt}
\caption{How this survey relates to existing surveys and position papers. ``LLM editor'': the work discusses a general-purpose language model choosing weight values directly. ``Lineage'': it connects WSP to hypernetworks, model editing, fast weight programmers and compiled transformers. ``RSI'': it discusses recursive self-improvement. ``Persistence'': it examines whether edits survive fine-tuning, quantization or perturbation. \checkmark\ = covered; $\circ$ = partial or thin; -- = absent. Checks are by full-text search and reading, as of November 25, 2026. The two Preprints.org surveys were read from Internet Archive copies of their full text, because the publisher's site refused automated access. Both WSP surveys connect WSP to hypernetworks, model editing and compiled Transformers but not to fast weight programmers, so we mark both as partial on lineage. The emecii survey is in Chinese; an English edition exists in an unmerged branch of its repository.}
\label{tab:surveys}
\begin{tabularx}{\linewidth}{>{\raggedright\arraybackslash}p{4.1cm} l >{\raggedright\arraybackslash}X c c c c}
\toprule
Survey & Date & Focus & LLM editor & Lineage & Persistence & RSI \\
\midrule
Weight space learning \citep{ax2603_10090} & 2026-03 & Weight geometry, embeddings, generation by trained models & -- & $\circ$ & -- & -- \\
Weight space as a modality \citep{ax2605_18632} & 2026-05 & Position: checkpoints as a generative modality & -- & $\circ$ & -- & -- \\
Model merging in the LLM era \citep{ax2603_09938} & 2026-03 & Merging, task arithmetic, sparsification & -- & $\circ$ & -- & -- \\
LLMs as optimizers \citep{ax2606_15577} & 2026-04 & Direct vs.\ tool-augmented LLM optimization & -- & -- & -- & -- \\
Safety in self-evolving agents \citep{ax2606_23075} & 2026-06 & Threats across agent modules and lifecycle & -- & -- & -- & $\circ$ \\
Self-improvement in agentic systems \citep{ax2607_13104} & 2026-07 & Model vs.\ scaffold updates & -- & $\circ$ & -- & \checkmark \\
RSI: bounded to autonomous research \citep{ax2607_07663} & 2026-07 & Loop closure, evaluators, auto-research & -- & -- & -- & \checkmark \\
The path to RSI agents \citep{liu2026pathrsi} & 2026-08 & Five-level RSI grading, co-improvement & -- & -- & -- & \checkmark \\
RSI in AI: a survey \citep{wu2026rsisurvey} & 2026-09 & Mechanism-level account of RSI & -- & -- & -- & \checkmark \\
WSP: methods, cost and limits \citep{emecii2026survey} & 2026-11 & Broad: WSP with editing, zeroth-order optimization, steering, architecture search & \checkmark & $\circ$ & $\circ$ & -- \\
WSP with LLMs: a survey \citep{turkoglu2026survey} & 2026-11 & Evidence audit of 70 studies (54 primary); WSP-12 checklist & \checkmark & $\circ$ & $\circ$ & $\circ$ \\
\midrule
\textbf{This survey} & 2026-11 & Definition, lineage, reliability agenda, RSI conditions, open problems & \checkmark & \checkmark & \checkmark & \checkmark \\
\bottomrule
\end{tabularx}
\end{table}


\paragraph{How to read this survey.} \Cref{sec:background} gives background. \Cref{sec:defs} defines the field and the taxonomy. \Cref{sec:methods} reviews methods by branch, and \cref{sec:apps} reviews applications by domain. \Cref{sec:eval} covers evaluation and reliability. \Cref{sec:rsi} discusses self-improvement, and \cref{sec:open} lists open problems. Appendix~\ref{app:table} lists every study with its codes, and Appendix~\ref{app:method} describes how we collected and read the literature. We intend this as a living survey: new versions will be dated, and the corpus, notes and generator scripts are released with the paper.
